% SPDX-FileCopyrightText: 2026 Will Cook
% SPDX-License-Identifier: CC-BY-4.0
% Open-source strategy paper; build with tectonic.
\documentclass[10pt]{article}
\usepackage{paper-house-style}
\usepackage{booktabs,array,float,placeins,multicol,graphicx}
\usepackage{tikz}
\usetikzlibrary{arrows.meta,positioning,calc,fit,shapes.geometric,matrix}
\usepackage[colorlinks=true,linkcolor=linkinternal,urlcolor=linkexternal,citecolor=linkinternal]{hyperref}
\hypersetup{
  bookmarksnumbered=true,bookmarksopen=true,bookmarksopenlevel=1,
  pdftitle={From Spare Compute to Cumulative Mathematics},
  pdfauthor={Will Cook},
  pdfsubject={An open-source strategy for agent-assisted mathematical research},
  pdflang={en-GB},
  pdfkeywords={open source, formal mathematics, Lean 4, AI agents, volunteer computing, research commons, attribution},
}
\newcommand{\repobase}{https://github.com/wcook04/plectis-erdos}
\newcommand{\repolink}[2]{\href{\repobase/blob/9b654f4cce44384f69d21fd2d0c51e62c4812b76/#1}{\texttt{#2}}}
\newcommand{\repoissue}[2]{\href{\repobase/issues/new?template=#1}{#2}}
\emergencystretch=1.7em

\runninghead{From spare compute to cumulative mathematics}
\title{From Spare Compute to Cumulative Mathematics}
\subtitle{An open-source strategy for agent-assisted research on hard problems}
\author{Will Cook}
\date{September 2026}

\begin{document}
\maketitle
\provenancenote{\emph{Authorship and AI use.} Every word of this manuscript was
generated by agents based on large language models operating within Will
Cook's private research system for artificial intelligence. Cook set the
objectives and acceptance criteria, reviewed and authorised the manuscript,
and is responsible for its contents. The agents are production tools, not
authors. This production disclosure supplies no independent verification;
Cook did not independently verify every mathematical claim.}
\fontsize{9.2pt}{11pt}\selectfont

\begin{abstract}
An open-source research environment lets contributors share the cost of
navigation, experiment records, formal proof checking, claim review, and
attribution.  Once these tools exist, someone with a mathematical idea or a
night of spare compute should be able to use them without rebuilding the
infrastructure.  This paper describes a strategy for doing so with
agent-assisted mathematical research.

Contributors may supply compute, mathematical direction, infrastructure
improvements, or review.  Useful returns include proofs, counterexamples,
finite computations, corrected statements, failed mechanisms, formal no-go
theorems, and research tools.  These can reduce the cost or improve the
direction of later work when their interpretation survives review.  Accepted
mathematics enters a versioned problem corpus; accepted infrastructure changes
alter how later work is produced and checked.  The record preserves the
origin of each contribution and the roles of those who supplied it.

The implementation is a public Lean repository organised around eight open
Erd\H{o}s problems, chosen to test the research process on difficult questions.
A fresh clone contains papers, formal source, explicit open boundaries, corpus
queries, validation programs, and structured return paths.  All eight problems
remain open.  The repository demonstrates research infrastructure; it supplies
no evidence that distributed agents outperform mathematicians or that more
compute will solve a named problem.  As of 31 August 2026, the author had
recorded neither a completed external cold-clone use nor an accepted external
contribution.  Contributor experience therefore remains untested outside
author-operated runs.

The account includes a production model, a contribution protocol with trust,
security, and attribution rules, and a proposed evaluation.  It compares the
project with volunteer computing, Polymath-style collaboration, formalisation
projects, and multi-agent proof systems.  Mathematics is the first domain
because its objects are digital and formal claims can be checked
incrementally.  Experimental science would require domain scientists,
physical laboratories, safety governance, and other sources of evidence.
\end{abstract}

\begin{paperresult}{What is implemented}
\textbf{Contribution.} A public contribution protocol accepts compute,
mathematical direction, infrastructure, and review as independently useful
inputs, while preserving evidence boundaries and role-specific credit.
\textbf{Implementation.} One fresh clone contains eight problem worlds,
formal source, explicit open obligations, validation programs, and structured
return paths.  \textbf{Limit.} All eight problems remain open, and no external
contribution has yet tested the proposed contributor experience.
\end{paperresult}

\section{The strategy}\label{sec:strategy}

This project began with one undergraduate working without an institutional
research team or a dedicated compute allocation.  Those circumstances explain
the interest in sharing infrastructure; they provide no validation of the
mathematics or the software.  The proposed research commons lets people share
an environment while contributing different kinds of work.

A mathematician may supply
a theorem, counterexample, reference, or correction.  A Lean contributor may
formalise or audit one statement.  A person with spare compute may run an
agent against a bounded public frontier; at present that means choosing a
bounded task, operating it through a runner of their own, and returning
evidence.  There is no turnkey client.  An infrastructure contributor may
improve navigation, experiments, validation, reproducibility, or the public
contribution path.  A reviewer may compare a formal statement with its
intended meaning or find an error in a returned result.
The credit record should describe each contribution on its own terms,
including work that leaves the problem open.

Contributors may use language models and agent harnesses throughout the work
and must disclose that use.  Model assistance alone says nothing about the
mathematical value of a return.  Acceptance concerns the attributable work:
an idea or direction that produces useful mathematics, a proof or
counterexample, a review, or an infrastructure change that improves later
research.  Each needs evidence and review appropriate to its claim.

The unresolved questions give the repository a continuing purpose.  A run
selects a subproblem, reads prior work, tests conjectures, records failed
routes, formalises stable steps, and explains what was established and what
remains open.  Even when it produces no solution, it can leave a record that
another researcher can inspect and use.

A \emph{continuous goal} keeps the endpoint fixed across runs while updating
the known results and remaining work.  The next agent reads the accumulated
theorems, computations, counterexamples, failed mechanisms, citations, and open
obligations before choosing an attempt.  Its return may be a theorem, a
narrower reduction, a refuted route, a source correction, or a better question.
The run is assessed by that return and its evidence.  Preserving it avoids
paying to reconstruct the same starting point.

\subsection*{Choosing what to automate}

The choice of task should state what someone hopes to understand and what
would change their next mathematical decision.  Recovering an earlier proof,
checking one calculation, comparing two mechanisms, or explaining a
counterexample may be the useful task.  A newly available solver is not by
itself a reason to expand the search.  When the purpose is learning, a hint or
a check of the reader's own attempt may serve it better than the completed
argument.  The contributor guide and coupled-goals skill state this
distinction; the reading guide offers an optional worked countermodel with
the hint and explanation separated.

The September 2026 Math and AI declaration identifies conceptual understanding,
the development of students and ideas, attribution, and human transmission as
purposes that a race to solve problems can undermine \cite{mathandai2026}.
For this project, that means choosing automation with the people doing the
work and respecting the terms on which a question was shared.  An open
discussion of a student's project or an unfinished approach is not a request
for an automated search to complete it.  A bounded collaboration can still
welcome substantial automation, but its purpose and expectations should be
clear before the search begins.

For a substantial search, the record should preserve the details needed to
interpret or reproduce it, including alternatives that explained a change of
direction.  Relevant details may include starting sources, tools, material
human interventions, effort and the stopping condition; estimates and unknowns
remain labelled.  Small corrections need no run history.  A reported success
rate must include unsuccessful attempts.  A timeout records a limit of that
run; a mathematical obstruction requires an argument at its stated scope.  These records can help compare
methods under declared conditions.  They do not establish a fixed boundary
between problems that AI can and cannot solve, a distinction raised by Tao in
his discussion of the changing difficulty landscape \cite{tao2026landscape}.

\subsection*{Discovery and stewardship}
\papersectiontarget{strategy-coupled-goals}

The proposed deployment separates two persistent roles.  A
\emph{discovery goal} stays close to one mathematical frontier: it reads the
corpus, compares attacks, uses computation to discriminate them, attempts
ordinary and formal proofs, and returns the smallest stable theorem,
counterexample, no-go, computation, or corrected boundary.  A
\emph{stewardship goal} stays close to the corpus as a whole: it compares the
new object with existing results, decides which declarations form one
mathematical family, reconciles the papers and assurance records, and returns
the next research question justified by that assessment.  The corresponding systems
contract is stated in the \papersectionlink{claim-faithful-publication-systems-paper.pdf}%
{systems-coupled-goals}{coupled-goals section of the systems paper}.

The two roles can be run by different people, agent harnesses, or machines.
A compute contributor may operate only the discovery side.  A mathematician
may suggest a direction, assess whether a result is nontrivial, or repair its
statement and exposition without paying for a long model run.  An expositor
can revise the paper, a Lean contributor can strengthen the formal interface,
and an agent builder can reduce the cost of a run.  The record credits the
work each supplied.

Stewardship determines where future effort goes as well as how a result is
presented.  It may discover that a new ``theorem'' is only a reformulation, that five declarations are one coherent
result family, that a short no-go eliminates an expensive research direction,
or that a paper still leads with a weaker theorem.  It then updates four
separate outputs: proof and evidence status, mathematical appraisal, paper
prominence, and the next allocation of compute or expert attention.  These
outputs may influence one another, but they are not one score and none can
promote an unproved claim.

The roles resume when there is new work to assess.  A stable mathematical
change calls for stewardship; a revised assessment, an affected result that
needs checking, or a sharper open boundary can justify another discovery run.
When the repository is unchanged, neither role needs an agent turn to report
that fact.

The public \repolink{skills/run-coupled-research-goals/SKILL.md}%
{run-coupled-research-goals skill} specifies how to run this sequence in a
cold clone.  It directs the mining and consequence-propagation jobs to use a
shared source pin and exchange committed objects and receipts.  One agent may
alternate between the roles; different people, models, subscriptions, or
machines may also perform them.  The decisions remain separate in either
arrangement.

The public clone must work without the private machine that produced the
initial corpus.  The private workbench is part of that production history and
grants no proof or publication authority to a return.  Contributors may use
any model runner or no model at all, provided their work follows the public
problem definition, evidence requirements, and review procedure.

The \papersectionlink{claim-faithful-publication-systems-paper.pdf}%
{systems-lifecycle}{systems paper} gives the complete claim-transition
lifecycle. The \papersectionlink{cold-clone-to-proof-receipt.pdf}%
{cold-clone-problem}{cold-clone paper} tests how a new agent reaches the public
mathematics and the validation boundary. This paper describes participation,
credit, and growth. The three accounts concern different parts of the same
prototype.

\subsection*{What one clone lets a contributor do}

A contributor does not need to understand the whole repository before doing
useful work.  From one clone, a person or agent can choose one of five first
actions: mine a bounded problem route; contribute mathematical direction
without paying for the compute that follows it; formalise or review one claim;
repair the research machinery; or propose another sourced problem world.  The
clone-local \repolink{skills/explain-public-system/SKILL.md}%
{explain-public-system skill} lets an agent read the public corpus and companion
papers on the newcomer's behalf, explain the claim boundaries at the requested
level, and point back to exact evidence.  The reader can therefore begin with
the ordinary request ``explain this repository to me'' rather than first
mastering its file layout.  The other skills select and run a frontier,
coordinate the coupled goals, install the same workflows in a compatible agent
harness, and describe the present multi-stage process for adding a problem.

The clone-local \repolink{scripts/agent_entry.py}{task-entry command} makes
that first choice inspectable. A request to ``maintain public infrastructure''
selects the maintenance workflow and its checks; a request to prove a remaining
implication selects bounded research and points to the relevant corpus support
and proof-plan queries. The output names the lane, a small set of documents to
read, the next commands, and the claim boundary before work begins. These routes
and their regression checks are implemented. Whether they reduce contributor
effort or improve distributed research remains unmeasured.

\begin{figure}[H]
\centering
\begin{tikzpicture}[
  action/.style={paper stage,text width=3.15cm,minimum height=15mm},
  arr/.style={paper arrow}
]
\node[action,fill=MidnightBlue!5] (mine)
  {{\scriptsize\sffamily\bfseries 1}\\[2pt]\textbf{Mine one route}\\[2pt]
   \scriptsize bounded theorem or obstruction};
\node[action,fill=MidnightBlue!5,right=6mm of mine] (direct)
  {{\scriptsize\sffamily\bfseries 2}\\[2pt]\textbf{Give direction}\\[2pt]
   \scriptsize hypothesis, example, or counterexample};
\node[action,fill=TealBlue!5,right=6mm of direct] (formal)
  {{\scriptsize\sffamily\bfseries 3}\\[2pt]\textbf{Check a claim}\\[2pt]
   \scriptsize formalise or review};
\node[action,fill=BurntOrange!6,below=6mm of $(mine.south)!0.5!(direct.south)$] (repair)
  {{\scriptsize\sffamily\bfseries 4}\\[2pt]\textbf{Repair the machinery}\\[2pt]
   \scriptsize navigation, validation, or exposition};
\node[action,fill=BurntOrange!6,below=6mm of $(direct.south)!0.5!(formal.south)$] (add)
  {{\scriptsize\sffamily\bfseries 5}\\[2pt]\textbf{Add a problem}\\[2pt]
   \scriptsize sourced world and boundary};
\node[paper box,below=8mm of $(repair.south)!0.5!(add.south)$,
      text width=10.8cm,fill=ForestGreen!6,align=center] (return)
      {{\scriptsize\sffamily\bfseries COMMON OUTPUT}\\[3pt]
       \textbf{Smallest inspectable delta}\\[2pt]
       \footnotesize source \(\cdot\) evidence \(\cdot\) limitation \(\cdot\) open statement};
\draw[arr] (mine.south) -- (repair.north);
\draw[arr] (direct.south) -- ($(repair.north)!0.5!(add.north)$);
\draw[arr] (formal.south) -- (add.north);
\draw[arr] (repair.south) -- (repair.south |- return.north);
\draw[arr] (add.south) -- (add.south |- return.north);
\end{tikzpicture}
\caption{Five useful first actions.  None requires an endpoint solution, and
none receives a stronger status than the evidence returned with it.}
\label{fig:first-actions}
\end{figure}

\section{A production model for mathematical progress}\label{sec:model}

Candidate production depends jointly on compute, model capability, human
direction, infrastructure, problem choice, and accumulated knowledge.
Mathematical progress additionally depends on validation and scarce human
review.  The notation below makes those dependencies explicit.  Let
\[
 C_t,\ A_t,\ H_t,\ I_t,\ P,\ K_t,\ V_t,\ R_t
\]
denote, at time $t$, available compute, agent capability, contributed human
mathematical direction, infrastructure quality, problem selection, accumulated
research knowledge, validation capacity, and review capacity.  The rate of
candidate generation may be written schematically as
\[
 G_t=G(C_t,A_t,H_t,I_t,P,K_t).
\tag{2.1}\label{eq:generation}
\]
The rate of \emph{reviewable mathematical progress} is a different quantity:
\[
 Y_t=Q(G_t,V_t,R_t),
\tag{2.2}\label{eq:accepted}
\]
where $Q$ includes formal checking, statement reconciliation, attribution,
and editorial selection.  Equations~\eqref{eq:generation} and
\eqref{eq:accepted} are a conceptual production model, not a fitted empirical
law.  They record complementarity and bottlenecks.  More compute can generate
more attempts without increasing the rate at which strong results are
validated or understood.  A better model can still revisit a closed route if
the corpus is difficult to navigate.  Expert insight can change the search
distribution without supplying any compute.  Review capacity can be the
binding constraint even when proof generation is cheap.

The open-source strategy adds two recursive updates:
\[
 \begin{aligned}
 K_{t+1}&=K_t+\Delta_t^{\rm theorem}
              +\Delta_t^{\rm counterexample}
              +\Delta_t^{\rm no\mbox{-}go}
              +\Delta_t^{\rm computation}
              +\Delta_t^{\rm correction},\\
 I_{t+1}&=I_t+\Delta_t^{\rm navigation}
              +\Delta_t^{\rm validation}
              +\Delta_t^{\rm reproducibility}
              +\Delta_t^{\rm governance}.
 \end{aligned}
\tag{2.3}\label{eq:feedback}
\]
Only reviewed returns with explicit scope enter these updates.  The first
line retains mathematical results for later research; the second records
changes intended to reduce cost or improve reliability.  A failed attempt may
supply a mathematical obstruction for $K_t$, or expose a recurring process
defect whose repair belongs in $I_t$.  Repairing the process does not change
the truth status of a conjecture.

\begin{figure}[H]
\centering
\begin{tikzpicture}[
  column/.style={paper box,align=center,text width=3.15cm,minimum height=5.0cm},
  receipt/.style={paper box,align=left,text width=5.05cm,minimum height=1.5cm,
                  font=\footnotesize},
  arr/.style={paper arrow}
]
\node[column,fill=MidnightBlue!5] (roles)
  {{\scriptsize\sffamily\bfseries OUTSIDE FUNCTIONS}\\[4pt]
   {\large\textbf{Inputs}}\\[6pt]
   \footnotesize mathematical direction\\[2pt]
   compute and agents\\[2pt]
   formalisation\\[2pt]
   infrastructure\\[2pt]
   review and exposition\\[2pt]
   new problems and funding};
\node[column,fill=TealBlue!5,right=8mm of roles] (engine)
  {{\scriptsize\sffamily\bfseries SHARED ENGINE}\\[4pt]
   {\large\textbf{Candidate production}}\\[6pt]
   \footnotesize \textbf{1} choose a frontier\\[3pt]
   \textbf{2} read the relevant corpus\\[3pt]
   \textbf{3} reason, compute, formalise\\[3pt]
   \textbf{4} return a candidate};
\node[column,fill=ForestGreen!6,right=8mm of engine] (result)
  {{\scriptsize\sffamily\bfseries PUBLIC RESULT}\\[4pt]
   {\large\textbf{Review funnel}}\\[6pt]
   \footnotesize replay and challenge\\[3pt]
   statement reconciliation\\[3pt]
   \rule{2.3cm}{.35pt}\\[5pt]
   \textbf{Accepted corpus}\\[2pt]
   public provenance};
\draw[arr] (roles) -- node[paper label,above]{contributions} (engine);
\draw[arr] (engine) -- node[paper label,above]{candidate} (result);

\node[receipt,below=7mm of $(roles.south)!0.5!(engine.south)$] (resource)
  {{\scriptsize\sffamily\bfseries RESOURCE RECEIPT}\\[2pt]
   human time; compute; model and harness; starting commit};
\node[receipt,right=7mm of resource] (lineage)
  {{\scriptsize\sffamily\bfseries CONTRIBUTION RECEIPT}\\[2pt]
   who supplied what; supporting evidence; later uses and repairs};
\node[paper box,below=6mm of $(resource.south)!0.5!(lineage.south)$,
      text width=10.9cm,align=center,fill=BurntOrange!6] (feedback)
  {Accepted mathematics enlarges the corpus; accepted architecture improves later runs.};
\draw[paper feedback] (feedback.west) to[out=180,in=-90] (roles.south);
\end{tikzpicture}
\caption{The conversion loop.  The labels on the left are functions, not
classes of people: one person may perform several, and several people may
perform one.  Search output crosses replay and review before it enters the
accepted corpus.  Mathematical returns enlarge the corpus; architecture
returns alter the engine.}
\label{fig:conversion}
\end{figure}

The model also gives a longitudinal use for the repository.  A fixed,
versioned problem world can be revisited as models, runners, and compute
budgets change.  Comparisons must disclose the starting commit, model,
scaffolding, number of attempts, compute, and review procedure.  Otherwise a
success says little about which variable changed.  A public failure record is
equally important: capability claims are badly distorted when successes are
announced and the number and cost of failed attempts are hidden
\cite{tao2026}.

\section{Why begin with hard mathematics}\label{sec:math}

Mathematics suits a first implementation because its objects, programs,
papers, and formal statements can all travel in one clone.
Many conjectures can be probed by exact computation.  A proof assistant can
check a stable formal step without waiting for an entire long argument to be
written informally and retrofitted later.  This does not make mathematical
meaning automatic: Lean checks the proposition in the source, not whether it
is the proposition the project intended to study.

Computation has a second role.  It can provide an artificial form of local
intuition to a system that is better at writing and running code than at
reproducing a mathematician's tacit judgement.  Exact experiments can expose
counterexamples, compare representations, and locate a narrow regularity.
The results guide which approach to attempt and which conjectures to pursue.
They leave the evidence class of an unbounded statement unchanged.  In Bayesian language,
an experiment can update which route deserves attention; it cannot assign a
formal posterior probability to a theorem or replace proof.

The eight current problems were selected as demanding research environments,
not as a claim that they were the eight most tractable or valuable open
problems.  Their papers expose materially different contribution points.

\begin{figure}[H]
\centering
\begin{tikzpicture}[
  card/.style={paper card,text width=5.05cm,minimum height=1.35cm}
]
\matrix[column sep=5mm,row sep=3mm] {
  \node[card,fill=MidnightBlue!4] {\textbf{68 \quad Factorial reciprocal series}\\[-1pt]
    \footnotesize carries and denominator escape}; &
  \node[card,fill=MidnightBlue!4] {\textbf{243 \quad Reciprocal-tail rigidity}\\[-1pt]
    \footnotesize discrete dynamics and recovery}; \\
  \node[card,fill=TealBlue!4] {\textbf{249 \quad Binary totient series}\\[-1pt]
    \footnotesize arithmetic tails and anti-concentration}; &
  \node[card,fill=TealBlue!4] {\textbf{251 \quad Prime-gap dyadic series}\\[-1pt]
    \footnotesize prime structure versus tail escape}; \\
  \node[card,fill=ForestGreen!5] {\textbf{257 \quad Mersenne-support subseries}\\[-1pt]
    \footnotesize achievement sets and rational targets}; &
  \node[card,fill=ForestGreen!5] {\textbf{269 \quad Running-lcm sums}\\[-1pt]
    \footnotesize finite-state structure and carries}; \\
  \node[card,fill=BurntOrange!5] {\textbf{1041 \quad Lemniscate geometry}\\[-1pt]
    \footnotesize local flow versus global path geometry}; &
  \node[card,fill=BurntOrange!5] {\textbf{1049 \quad Rational-base series}\\[-1pt]
    \footnotesize Lambert approximation kernels and obstructions}; \\
};
\end{tikzpicture}
\caption{Eight distinct research environments.  The results map and the
problem papers state their exact frontiers; this strategy paper does not
reproduce the specialist vocabulary needed to attack them.}
\label{fig:problem-worlds}
\end{figure}

Every row admits more than theorem proving.  A mathematician may sharpen a
hypothesis or supply a counterexample.  A programmer may run a discriminating
exact calculation.  A formaliser may turn a paper deduction into a checked
declaration or find that the two statements disagree.  A literature reader
may locate a theorem that closes or invalidates a route.  Each return must
state what was established and its limitation; the problem number supplies
no additional evidence of value.

\section{The public research object}\label{sec:object}

The unit of participation is a \emph{problem world}.  A problem world includes
the endpoint question, current public status, principal theorems, formal
source, experiments, known counterexamples, failed mechanisms, source
literature, and exact open obligations.  A bounded task selects the part of
that world needed for one question.

The research was initially developed inside a private workbench.  The public
release carries both the mathematical corpus and the workflows needed to
inspect, validate, and return a contribution.  A contributor uses the public
checkout as the working environment; access to the private workbench is not
part of the contribution protocol, and the workbench grants no proof or
publication authority to a returned result.

Figure~\ref{fig:authority-ladder} separates the kinds of evidence and review
available in the public project.

\begin{figure}[H]
\centering
\begin{tikzpicture}[
  rung/.style={paper stage,text width=3.2cm,minimum height=2.35cm},
  wide/.style={rung,text width=4.95cm},
  arr/.style={paper arrow}
]
\node[rung,fill=MidnightBlue!5] (exp)
  {{\scriptsize\sffamily\bfseries 1 \; FINITE EVIDENCE}\\[3pt]
   \textbf{Experiment}\\[3pt]\footnotesize establishes a finite output\\[2pt]
   \textit{Scope: no unbounded theorem}};
\node[rung,fill=TealBlue!5,right=6mm of exp] (lean)
  {{\scriptsize\sffamily\bfseries 2 \; FORMAL EVIDENCE}\\[3pt]
   \textbf{Lean kernel}\\[3pt]\footnotesize accepts the exact proposition\\[2pt]
   \textit{Scope: intended meaning is separate}};
\node[rung,fill=ForestGreen!5,right=6mm of lean] (match)
  {{\scriptsize\sffamily\bfseries 3 \; ALIGNMENT}\\[3pt]
   \textbf{Statement review}\\[3pt]\footnotesize compares formal and informal claims\\[2pt]
   \textit{Scope: no novelty judgement}};
\node[wide,fill=BurntOrange!6,below=8mm of match] (expert)
  {{\scriptsize\sffamily\bfseries 4 \; SCHOLARLY REVIEW}\\[3pt]
   \textbf{Expert assessment}\\[3pt]\footnotesize soundness, relevance, and relation to prior work};
\node[wide,fill=Violet!5,left=7mm of expert] (community)
  {{\scriptsize\sffamily\bfseries 5 \; EXTERNAL OUTCOME}\\[3pt]
   \textbf{Community uptake}\\[3pt]\footnotesize develops through criticism, use, and reuse};
\draw[arr] (exp) -- (lean);
\draw[arr] (lean) -- (match);
\draw[arr] (match) -- (expert);
\draw[arr] (expert) -- (community);
\end{tikzpicture}
\caption{Evidence and acceptance remain separate.  A candidate need not pass
through every box in a single line, but no earlier box inherits the authority
of a later one.}
\label{fig:authority-ladder}
\end{figure}

The project builds on established mathematical and computational practices.  Erd\H{o}s Problems supplies questions, sources, status,
and a problem community \cite{erdosproblems}.  Lean supplies the formal
language and kernel \cite{lean4}.  DeepMind's
\texttt{formal-conjectures} supplies a public one-problem-file model with
metadata, snapshots, issues, and pull requests \cite{formalconjectures}.
Comparator and Palomar supply exact-interface checking, permitted-axiom
inspection, and a durable formal record with a deliberately limited editorial
claim \cite{palomar}.  Polymath supplies a social precedent for publishing
small, tentative, and negative contributions \cite{polymath}.  BOINC and GIMPS
supply the volunteer-compute precedent and make visible why executable work,
independent checks, and legible credit matter \cite{boinc,gimps}.  The present
repository wraps these practices into a route that a newcomer can enter
without first rebuilding each component.

The same sources can serve a short public primer, a specialist paper, a
detailed proof account, Lean declarations, an agent explanation, and
machine-readable queries.  Every account must link to its source claims and
evidence and preserve their scope, however briefly or informally it explains
them.

An external runner begins at \repolink{AGENTS.md}{the compact agent
entry}.  It can inspect all problem frontiers, choose a bounded question, read
the relevant paper and source neighbourhood, run experiments or edit formal
code, validate the result, and prepare a return.  None of these steps requires
access to the private workbench.  A contributor is free to replace the agent,
the scheduler, or the entire search policy while retaining the public evidence
boundary.

\section{The contribution protocol}\label{sec:protocol}
\papersectiontarget{strategy-protocol}

To propose a patch, a contributor forks or clones the repository, creates a
branch, commits a focused change, pushes the branch to a public fork, and
opens a pull request.  A person who has a result but no patch
can open a structured research-progress issue.  A person with an
infrastructure proposal can use the parallel architecture-proposal path.

The clone-local \texttt{submit-pull-request} skill gives agents the same
procedure. It groups separable mathematical or infrastructure changes into
coherent commits, runs the relevant checks, and prepares the repository's
pull-request template. A commit may span several files when they form one
change. Pushing to a fork and opening the pull
request are external actions and occur only after the contributor authorises
them.

Upstream work may continue while the contributor's branch remains open. The
recorded starting commit supplies the common ancestor from which Git recovers
the original delta; the contributor does not need a copy of the maintainer's
private or current working tree. Review therefore has two passes. First,
reproduce the change in its original public context. Second, reconcile the
accepted substance with current main and rerun the present checks. A material
conflict resolution is a new, separately credited integration change. It does
not erase or retrospectively rewrite the original contribution.

The complete protocol has nine steps.

\begin{enumerate}
\item \textbf{Clone.}  Record the exact public commit and environment.
\item \textbf{Orient.}  Read the claim boundary and the paper governing the
      selected problem or architecture area.
\item \textbf{Select.}  Name one bounded frontier, expected evidence, and a
      stop condition.  Working on all problems is allowed, but each returned
      result remains independently scoped.
\item \textbf{Work.}  Use any mixture of reasoning, source research,
      computation, Lean, and human mathematical judgement.  Preserve useful
      negative results instead of rewriting the run as a linear success.
\item \textbf{Validate.}  Run the checker appropriate to the claim.  A
      computation receives a computation receipt; a formal theorem receives a
      Lean receipt; neither receives semantic or publication authority from
      that fact alone.
\item \textbf{Propagate.}  Enumerate the plausible Lean, claim, paper,
      computation, route, validation, and contributor consumers.  Update each
      one, verify it unchanged, defer it with a re-entry condition, or explain
      why it is not a consequence.  Repeat this pass after old-branch work is
      reconciled with current main.
\item \textbf{Return.}  Submit source, evidence, starting commit, tool
      disclosure, collaborators, requested contribution roles, and the exact
      stronger statement that remains unproved.
\item \textbf{Review and adopt.}  Reviewers reproduce the evidence, reconcile
      formal and informal statements, decide scope, and either accept, request
      revision, reject, or preserve the return as a scoped negative result.
\item \textbf{Publish and credit.}  An accepted generation is linked to the
      paths and commit that assimilated it.  Later corrections append lineage;
      they do not erase the earlier contributor.
\end{enumerate}

A non-specialist compute contributor needs an intelligible task and a way
to return evidence.  They are not asked to certify a proof or invent a
mathematically meaningful objective unaided.  The default work unit is a
packet prepared from an authored mathematical route: an exact open statement,
relevant sources, a permitted experiment or proof obligation, an expected
evidence class, and a stopping condition.  Qualified mathematical review should precede large-scale
distribution; it is not claimed for every present route.  The contributor and
agent execute the packet and return evidence.
Lean checks a formal proposition.  A qualified reviewer checks whether that
proposition says what was intended and whether the return is relevant.  The
return remains a candidate until the applicable gates have been crossed.

The mathematician who supplies an input and the mathematician who verifies a
return are also distinct roles.  The first may contribute a conjectural
mechanism, a counterexample pattern, or simply a well-posed question that
causes the system to produce useful mathematics.  The second inspects the
resulting statement, argument, prior art, and significance.  One person may
perform both roles, but the provenance record must not let an input certify
its own output merely because it came from a mathematician.

The protocol accepts two kinds of contribution.  The \emph{mathematics track}
accepts proofs, reductions, computations, counterexamples, corrections,
formalisation, and reproducible failed routes.  The \emph{architecture track}
accepts improvements to agent workflow, navigation, validation,
reproducibility, public experience, governance, and tooling.  An architecture
return never needs a fictitious problem number, and it cannot request
promotion of a mathematical claim.

A pull request proposes a change; review records whether it is adopted, and
a release records an immutable public version.  These separate records permit
credit for a useful counterexample, correction, or design even when no theorem
is merged.

\section{From an agent claim to mathematical attention}\label{sec:attention}

An agent saying ``solution found'' is an event in a run log, not a project
claim.  An informal derivation may contain an unnoticed gap.  A Lean theorem
may be completely kernel-correct while formalising the wrong statement,
assuming an inadmissible axiom, or proving a result too weak to settle the
problem.  Formal verification therefore strengthens one layer of evidence;
it does not collapse correctness, intended meaning, novelty, significance,
and community acceptance into a single status.

The proposed review sequence reserves specialist attention for returns that
have passed the earlier checks:

\begin{enumerate}
\item \textbf{Candidate.}  A person or agent returns a scoped argument,
      computation, formal declaration, counterexample, or no-go.
\item \textbf{Replay.}  Automated checks reproduce the code and attach the
      correct evidence class without promoting its claim.
\item \textbf{Internal challenge.}  Independent agents or contributors try
      to break the statement, locate stronger prior art, vary the experiment,
      and reconcile formal and informal formulations.
\item \textbf{Maintainer triage.}  A human checks whether the bundle is
      coherent and relevant, and whether the earlier checks justify asking
      a specialist to spend time on it.
\item \textbf{External record.}  A mature Lean result may be packaged through
      Comparator and submitted to Palomar; a mathematical account relevant to
      a listed problem may be placed before the Erd\H{o}s Problems community
      and the appropriate research audience.
\item \textbf{Acceptance.}  Mathematicians inspect, explain, criticise, reuse,
      publish, or reject the work over time.  Broad mathematical acceptance is
      a community judgement that the repository maintainer cannot grant.
\end{enumerate}

Palomar mechanically verifies a pinned challenge--solution interface using
Comparator and kernel checking, records structured disclosure, and applies a
documented automated editorial filter.  It explicitly does not endorse the
result, establish novelty, or provide human expert review \cite{palomar}.
Likewise, the Erd\H{o}s Problems forum asks long proofs or partial proofs to be
linked as external documents and filters implausibly incomplete claims; it is
a route to relevant attention, not a substitute for that community's
judgement \cite{erdosproblems}.

Several contributors can make a return easier to assess through independent
replay, adversarial review, repairs to different failure modes, a stable
formal interface, or exposition that a specialist can audit.  Their number
alone provides no mathematical evidence.  The work they record can justify
specialist attention without strengthening the theorem's status.  This gives
a contributor without institutional connections concrete reasons to ask for
review, while sparing the expert the raw agent transcripts.

\section{Progress, provenance, and credit}\label{sec:credit}
\papersectiontarget{strategy-credit}

Each return states what kind of evidence it supplies.  The accepted forms
include:

\begin{itemize}
\item an unconditional theorem or a formal proof of an existing statement;
\item a counterexample or corrected hypothesis;
\item a no-go theorem excluding a defined strategy class;
\item a bounded computation with code, range, and interpretation;
\item a conditional reduction with its antecedent still visible;
\item a reproducible failed route that prevents duplicated search;
\item a literature correction or exact source locator;
\item an architecture proposal or validated infrastructure change; or
\item exposition or review that changes what another person can understand
      and check.
\end{itemize}

A single leaderboard would obscure the differences between compute volume,
mathematical depth, software quality, correction work, exposition, and review.
The public recognition view can list contributors with their accepted
artefacts, evidence classes, problem or infrastructure area, and roles.  The current receipt schema supports the
fourteen CRediT roles as an optional vocabulary, including conceptualisation,
methodology, software, validation, investigation, and writing
\cite{credit}.  CRediT describes work; it does not decide authorship.

Adopting a contribution does not give the project ownership of its underlying
idea.  The public record should name its originator and identify the work
that was retained.  Authorship and formal publication still require separate
judgements.  Recording roles preserves contributions that would disappear
under sole solver credit or an undifferentiated collective name.

The same rule applies when a local result is carried into a broader library.
The ledger should distinguish the person or run that produced the local
observation from the expert who recognised its natural generality, reconciled
it with prior art, designed the library interface, rewrote or formalised the
proof, and stewarded review.  If that expert turns a problem-specific result
into a Mathlib-quality contribution, the project does not demand ownership of
the resulting pull request or paper.  It asks for a reasonable provenance
link back to the route from which the idea arose, while crediting the expert
fully for the mathematical judgement and upstream work they actually
performed.

Credit for partial and negative work gives an outside contributor a reason
to spend thought or compute on an unsolved problem.  The same record makes
later correction, reproduction, and literature attribution possible.

A contributor outside a university, research institute, or AI laboratory can
point to the accepted artefact.  Its receipt identifies their idea or
implementation, the public files containing it, the evidence reviewers
checked, their roles, and the remaining boundary.  The record documents that
contribution publicly.  It confers no degree, peer-review status, employment
recommendation, or guarantee of recognition by another institution.

\subsection*{Record now, trace consequences later}

The first ledger should remain descriptive.  A contribution receipt records
who supplied which object, from which public state, under which role, with
which evidence and limitation.  It need not guess the contribution's eventual
importance.  Later accepted work can add reviewed lineage edges such as
\emph{uses}, \emph{enables}, \emph{repairs}, \emph{formalises}, \emph{refutes},
or \emph{explains}.  A consequence view can then show which later results,
tools, or expositions depend on an earlier contribution.

Before returning a result, the contributor follows the clone-local
\texttt{propagate-research-consequences} skill to inspect its plausible uses.
Each affected statement, file, or process is updated, verified unchanged,
deferred with a reason, or excluded from scope.  Later work may reveal further
consequences; known uses still need checking at the time of return.

These links provide a revisable account of later uses, without assigning an
objective impact score.  A later event does not establish causation, and an
authored dependency claim can be mistaken.  A small observation may have many
consequences; a large compute run may have none beyond its negative result.
Both keep their original receipt, and readers must be able to inspect and
correct the subsequent links.

Any financial reward scheme would require a separate prospective policy for
eligibility, conflicts, funding, tax, dispute, and the treatment of old
receipts.  Existing attribution creates no financial entitlement.  A mining
run can still disclose the model, provider, harness, compute donor, starting
commit, and any human direction as distinct roles.

\section{Running contributed code safely}\label{sec:security}
\papersectiontarget{strategy-security}

BOINC packages scientific jobs for heterogeneous consumer devices, and BOINC
Central explicitly aims to
make volunteer computing available to scientists without the resources to
operate their own project \cite{boinc}.  GIMPS shows a mathematical version:
volunteers run a common search program, independent machines confirm a prime,
and discovery credit includes the compute donor, software authors, server
operator, and wider volunteer effort \cite{gimps}.

Agent-assisted research is less uniform than either example.  Most tasks are
not interchangeable work units with a predetermined verifier.  A runner may
change code, propose a conjecture, or misread the problem.  The public system
therefore distributes \emph{search} while keeping authority local to each
evidence type.  Mathematicians and formalisation contributors design or review
the routes; volunteers may execute them with Claude Code, Codex, Cursor,
Antigravity, OpenCode, another open or proprietary harness, or no agent at all.  The
shared task packet, evidence boundary, and return record permit different
runners.  Returned code is untrusted until reviewed.  Expensive or privileged
continuous-integration jobs must not execute fork code with
repository secrets.  In particular, a GitHub \texttt{pull\_request\_target}
workflow must not check out and execute an untrusted fork; GitHub documents
that combination as a route to secret leakage and repository compromise
\cite{githubsecurity}.  Initial checks should run with read-only permissions,
no secrets, bounded resources, and explicit artefact retention.  Promotion to
trusted infrastructure is a later review decision.

A useful compute return records at least the source commit, model and runner,
tool disclosure, prompt or task packet, wall-clock time, human time, token or
subscription use, and hardware budget where available, together with commands,
changed paths, outputs, validation receipts, failures, and the stopping rule.
A mathematician's time, an API bill, and donated hardware cannot be reduced
to one exchange rate.  Reporting them separately still permits interpretable
comparisons and makes the cost of the surrounding assistance visible.

The first public mode may remain deliberately simple: contributors clone the
repository and run their own agents locally.  A later volunteer-compute layer
could distribute signed, immutable task packets and receive result bundles
without granting write access.  It should be built only after task identity,
sandboxing, deduplication, resource budgets, result replay, and abuse handling
have explicit owners.

\section{Precedents}\label{sec:prior}

The strategy combines practices from Erd\H{o}s Problems, Lean and mathlib,
Comparator and Palomar, DeepMind's \texttt{formal-conjectures}, Polymath,
BOINC, GIMPS, and recent multi-agent formalisation systems.  Its intended
contribution is to make these practices available through one clone.

\paragraph{Distributed compute.}
BOINC and GIMPS show that members of the public will donate hardware to a
scientific objective when the client is easy to run, the work is visible, and
credit is legible \cite{boinc,gimps}.  Their tasks are much more mechanically
uniform than open-ended proof research.

\paragraph{Distributed mathematical insight.}
Polymath projects invite participants at different mathematical levels to
share small observations as they occur.  Their rules explicitly welcome
tentative and negative insights, provided they are made clear enough for
others to absorb, and treat the project as collaboration rather than a race
\cite{polymath}.  Polymath supplies a social protocol; it does not supply a
formal proof, computation, and agent-return substrate for every comment.

\paragraph{Distributed formalisation.}
Mathlib uses ordinary fork-and-pull-request practice, human review, and
source-level attribution.  The Carleson formalisation used a public blueprint
and many claimable lemma-sized tasks, with formalisation feeding corrections
back into the informal plan \cite{mathlib,carleson}.  Google DeepMind's
\texttt{formal-conjectures} repository similarly uses one-problem files,
issue assignment, pull requests, metadata, and stable snapshots for a growing
Lean benchmark \cite{formalconjectures}.  These projects show how a large
formalisation can be divided into tasks for public contributors.

\paragraph{Multi-agent formal research.}
Agent Hunt studies bounties, locks, guarded ownership, and collaborative
agents for autoformalisation \cite{agenthunt}.  Lean Atlas uses formal
dependency information to reduce the declarations a person must inspect for
semantic verification \cite{leanatlas}.  These mechanisms address search,
coordination, and review focus.  The
additional proposal here is a public adoption procedure that records the
provenance of both mathematics and improvements to the research process.

\paragraph{Proof abundance.}
Tao separates problem solving into generation, verification, exposition,
digestion and acceptance, and canonicalisation, and argues that proof
abundance will create bottlenecks between these stages \cite{tao2026}.  The
strategy invites contributions at each of these stages, including review
that prevents an overclaim and exposition that explains a hard step.  Tao's
rule of thumb, that a result
whose authors cannot give a clear, correct, properly attributed expert-level
account of it is incomplete even when formally verified, sets the standard
for exposition here.  A record designed to help a reader learn a result does
not establish that its authors have supplied such an account.  That must be
shown for each result.  At the time of writing, no selected result has a
recorded expert account of this kind.  Section~\ref{sec:evaluation} proposes
testing whether the record helps people produce one.

A digestive-system analogy helps distinguish the preparation from its use:
breaking down food is different from absorbing its nutrients.  The AI system
performs \emph{pre-digestion} when it organises sources, expands a derivation,
prepares examples, or connects a statement to its proof.  People can do this
preparatory work too.  Mathematical digestion occurs when mathematicians work
through the argument, understand why its assumptions and difficult steps are
needed, and can explain and use its ideas.  A readable paper shows that
material has been prepared; it does not demonstrate human understanding.
We use this analogy to express Tao's separation of verification, exposition,
and community digestion and acceptance
\cite[Sec.~6, pp.~6--8; Sec.~8, p.~11]{tao2026}.  It describes the
preparation for review.  Claims still require verification, and human
understanding must be demonstrated through what people can explain and use.

The implementation brings volunteer compute, small shared insights, formal
task decomposition, problem worlds, recorded negative results, untrusted
returns, credit by role, and human review into one open corpus.  This
combination is the proposed contribution.  The paper makes no universal
priority claim and reports no controlled comparison showing an increased
discovery rate.

\section{Adding another problem world}\label{sec:add-problem}
\papersectiontarget{strategy-add-problem}

To extend the repository beyond the present eight problems, each new problem
needs a canonical question and primary source; a dated account of its current status; an explicit
public claim boundary; a readable primer or problem note; formal statements
with an informal-faithfulness account where formalisation is appropriate;
known literature, computations, counterexamples, and failed routes; exact open
contribution points; navigation and validation routes; attribution; and a
maintainer or review path.

The public skill for adding a problem separates three states.  A proposal may
begin as a sourced issue.  An incubating formal lane may add checked source
under a problem-owned namespace while stating that it is not yet a reviewed
public claim.  A fully indexed world joins the problem registry, paper corpus,
query routes, return schema, cold-start inventory, and relevant external
crosswalks.  These are distinct transitions because a checked proposition, a
published paper, a reviewed claim, and a Comparator selection are different
facts.

After any of these transitions, consequence propagation inspects the problem
registry, paper corpus, query routes, cold-start inventory, validation roster,
return schema, and external crosswalks. A new problem is complete at its stated
entry level only when the contributor has recorded what was done about each
affected use.

The last transition is not yet a one-command public operation.  The current
paper-corpus fan-in has a maintainer-owned stage, several checks enumerate the
present roster, and the bounded problem index is already close to its size
limit.  The immediate infrastructure work is therefore to derive rosters from
one public authority, admit an explicit incubating state, split verbose detail
out of the bounded index, and make an absent external crosswalk record a normal
state rather than an error.  The add-problem skill identifies these dependencies so that a contributor
can account for them during integration.

The design aims to accommodate increasingly deep problem records, but its
scaling has not been measured.  A useful test asks whether a new agent can
retrieve what it needs without reading every paper or rediscovering a known
failure as the number and depth of problems grow.  That test should be repeated whenever the roster or
navigation machinery changes.

\section{Participation and growth}\label{sec:growth}

The project should present several independent reasons to clone the
repository.

\begin{itemize}
\item A \textbf{mathematician} can see eight exact frontiers and contribute a
      lemma, construction, counterexample, reference, or critique without
      learning the private orchestration system.
\item A \textbf{Lean user} can formalise a paper deduction, audit a statement,
      simplify a proof, or improve the checked interface.
\item A \textbf{compute hobbyist} can run a local agent against a bounded route
      in a real open problem and return reproducible evidence even when it is
      negative.  The contributor is not asked to certify a proof claim;
      mathematical review remains downstream.
\item An \textbf{agent builder} can compare runners on a stable problem world
      while disclosing compute, attempts, and starting state.
\item An \textbf{infrastructure contributor} can improve the conversion from
      all other inputs to useful mathematical output.
\item A \textbf{reviewer or expositor} can reconcile meaning, rank results,
      correct attribution, or turn a checked proof into recoverable
      understanding.
\end{itemize}

The first screen of the README should give the basic invitation: clone the project, point a local agent at the compact entry file, and
ask it to choose a bounded frontier that matches the available tools.  The
same screen must say that all eight problems remain open, that a useful return
need not solve one, and that infrastructure contributions receive credit.

The contribution procedure should work before the project is publicised
through Hacker News, a research talk, or a model-community post.  The claim
would then be specific: an outside contributor can clone the corpus, find an
exact open question, run or improve the tools, submit a result with its
evidence, and find accepted work in the public record.  Stars, clone counts,
and agent-hours measure attention or activity, without establishing
mathematical progress.

\section{Evaluation}\label{sec:evaluation}

Evaluation should examine what happens between choosing a task, producing a
return, reviewing it, and using the result.  Initial measures include:

\begin{itemize}
\item time from a cold clone to a correctly scoped first task;
\item fraction of returned bundles that can be replayed at the stated commit;
\item fraction requiring statement or evidence-class correction;
\item time from return to first substantive review and to adoption decision;
\item repeated-dead-end rate before and after a no-go enters the corpus;
\item reviewer time per accepted result family;
\item contribution diversity across mathematics, computation, software,
      validation, exposition, and review;
\item correction lineage completeness and attribution retention; and
\item change in useful output at controlled compute when navigation or another
      infrastructure component changes.
\end{itemize}

Counts need mathematical interpretation.  One accepted counterexample may be
more useful than hundreds of generated lemmas.  Generated certificate shards
must not be counted as independent discoveries.  A theorem accepted by Lean
but rejected on intended meaning is not a successful public transition.
Likewise, a strong negative result can improve the corpus even though the
endpoint problem remains open.  A retained attempt is not yet knowledge: a
failed route may rest on a mistaken diagnosis, duplicate known work, or add
nothing actionable.  It counts only through a justified interpretation, such as
a counterexample, a scoped obstruction, a reproducible stopping point, or a
specific correction to the next task, and the record separates what happened
from what was learned.

Longitudinal model comparisons should use immutable problem snapshots and
held-out variants where possible.  Public hard problems are susceptible to
training-data contamination, and the repository itself will become more
informative over time.  A later run may therefore use both a stronger model
and a more informative corpus.  The study should measure these factors
separately before attributing later success to the model.

This paper reports the implementation and protocol and proposes measures
for a later study.  It supplies no such evaluation.

A small first study should diagnose difficulties in the contribution path;
it would not estimate their frequency across a population.  An outside researcher receives one problem world and is asked to
recover its strongest established statement, explain the main idea, state what
remains unproved, distinguish a genuine obstruction from an unsuccessful
attempt, and then undertake one bounded task and return it in a form the
maintainer can assess without substantial reconstruction.  Three presentations
of the same material should be compared: the source, papers, and records in an
ordinary repository; an information-equivalent static briefing written with
care; and the navigation and record surfaces described here.  The second
condition is what separates the effect of the mechanisms from the effect of a
better introductory paragraph.  Failures are recorded at their own level: a
command that could not be run, a source that could not be located, an
interpretation that was corrected, or a result judged uninteresting are four
different findings.

That study should also ask what the reader can do after the session: reconstruct
the decisive step without copying it, explain why a tempting alternative fails,
locate the prior idea, or formulate a nearby question with a mathematical
reason for asking it.  Record the actual explanation, correction or follow-up
use and the help the reader needed.  A completed checklist, shorter reading
time or a model's judgement that the prose is clear cannot substitute for
those observations.  Nor should these activities become a compulsory score
for exploratory work.  They are proposed ways to learn whether the record
helps people think; this paper reports no such reader outcomes.

\section{Learning from every run}\label{sec:learning}

The research memory has three levels.  Problem-local memory records the exact
objects that change the mathematical search: proved lemmas, counterexamples,
finite data, failed mechanisms, and unresolved obligations.  Reusable
mathematical memory records results that have survived a separate
generalisation and integration pass.  General process memory records lessons
that transfer: a navigation failure, an experiment pattern, a formalisation
pitfall, a validation gap, or a review rule.  A local theorem does not enter
the second level merely because its variables have been renamed.

A continuing discovery run should:

\begin{enumerate}
\item re-read the current frontier and the relevant corpus neighbourhood;
\item check whether the proposed route, calculation, or failure is already
      recorded;
\item compare distinct attacks and run the cheapest exact computation that can
      distinguish or refute them;
\item carry out the warranted analytic and incremental Lean work;
\item attack the strongest candidate for dropped hypotheses, counterexamples,
      and stronger prior art;
\item return the smallest evidence-bearing delta with its starting state and
      surviving boundary; and
\item add problem-local knowledge while proposing reusable process lessons
      separately for review.
\end{enumerate}

Stewardship should then:

\begin{enumerate}
\item wake on a stable theorem, counterexample, computation, authority change,
      paper correction, or review outcome rather than on a timer;
\item rebuild the relevant candidate universe from source authority, including
      stronger results absent from the present paper or Comparator roster;
\item group declarations into coherent mathematical families and distinguish
      a new mechanism from a renamed residual, routine corollary, or duplicate;
\item decide proof authority, mathematical significance, exposition order,
      and future work allocation separately;
\item reconcile the paper, Comparator interface, Palomar disposition, claim
      boundary, navigation, and contributor lineage wherever the result has a
      real consequence; and
\item return a source-pinned frontier update: the strongest surviving result,
      its hard step, the exact boundary, and the next discriminating question.
\end{enumerate}

A changed assessment of the mathematics can warrant a different paper order.
The strongest exact results and mechanisms should lead, while routine
supporting material remains available with less space.  Comparator checks
exact interfaces and Palomar provides an external registration route.
Stewardship may prepare and prioritise submissions to them; it cannot confer
novelty, acceptance, or canonical status.

Subagents can divide literature reading, computation, proof search,
formalisation, and adversarial review when their questions and evidence remain
independent.  The integrating agent must read and verify their returns, and
each lane keeps its own starting state and stop condition.  Every attempt
needs the same evidence regardless of how many agents contributed.

The system calls the transfer of a local process lesson to a reusable rule
\emph{up-propagation}.  A proposal states the evidence, comparable cases,
intended scope, and cases to which it should not apply.  Review determines
whether it warrants a change to a general skill or route.  Computational
reasoning can likewise become a specialised skill with reusable experiment
forms, but every new experiment still records its finite domain and
interpretation.

\subsection{The local-to-general track}\label{sec:local-to-general}
\papersectiontarget{strategy-local-to-general}

A theorem developed while studying a hard problem may be useful elsewhere,
even if that problem stays open.  Preparing it for broader use requires a
separate mathematical review:

\begin{enumerate}
\item identify the exact local result and retain its problem-specific proof and
      provenance;
\item separate the hypotheses genuinely used from the coordinates and names of
      the motivating problem;
\item search the literature and the target library for an existing theorem,
      collision, preferred abstraction, and likely downstream consumers;
\item state the natural general result, with the local theorem exhibited as a
      transparent special case rather than hidden by the abstraction;
\item falsify proposed weakenings, check the general proof independently, and
      test at least one other real use; and
\item ask a qualified mathematician and Lean contributor to decide whether the
      result is important, well-shaped, documented, maintainable, and ready for
      an upstream discussion.
\end{enumerate}

\begin{figure}[H]
\centering
\begin{tikzpicture}[
  node distance=6mm and 9mm,
  box/.style={paper stage,text width=3.3cm,minimum height=1.65cm},
  gate/.style={box,fill=BurntOrange!7,text width=3.7cm},
  arr/.style={paper arrow}
]
\node[box,fill=MidnightBlue!5] (world)
  {{\scriptsize\sffamily\bfseries LOCAL RESULT}\\[3pt]
   \textbf{Hard-problem world}\\[2pt]\scriptsize proof, no-go, or computation};
\node[gate,right=of world] (digest)
  {{\scriptsize\sffamily\bfseries GATE 1}\\[3pt]
   \textbf{Mathematical digestion}\\[2pt]\scriptsize hypotheses, prior art, falsifiers};
\node[box,fill=TealBlue!5,right=of digest] (generic)
  {{\scriptsize\sffamily\bfseries CANDIDATE}\\[3pt]
   \textbf{Reusable theorem}\\[2pt]\scriptsize natural statement and consumers};
\node[gate,below=of generic] (expert)
  {{\scriptsize\sffamily\bfseries GATE 2}\\[3pt]
   \textbf{Expert stewardship}\\[2pt]\scriptsize meaning, fit, maintenance, credit};
\node[box,fill=ForestGreen!6,left=of expert] (library)
  {{\scriptsize\sffamily\bfseries SHARED OUTPUT}\\[3pt]
   \textbf{Library result}\\[2pt]\scriptsize reused, or declined with reasons};
\draw[arr] (world) -- (digest);
\draw[arr] (digest) -- (generic);
\draw[arr] (generic) -- (expert);
\draw[arr] (expert) -- (library);
\end{tikzpicture}
\caption{A local result becomes reusable only through a second mathematical
and social gate.  Formal correctness is necessary where Lean applies, but it
does not decide generality, library fit, or stewardship.}
\label{fig:local-to-general}
\end{figure}

Mathlib provides an external model for this final stage.  Its contribution
guide welcomes useful
contributions and supplies a public fork-and-pull-request route; it does not
reserve contribution to people with institutional credentials.  It also sets
high standards for generality, integration, maintainability, style, and
documentation, and its current AI policy rejects low-quality unsupervised
agent submissions while requiring AI use to be disclosed \cite{mathlib}.
A compute donor should therefore not be directed to submit unreviewed agent
lemmas to Mathlib.  A subject and Lean expert must understand the candidate,
decide whether it belongs, bring it to community standards, and take
responsibility for the discussion.  Credit follows the roles they performed.
The originating problem remains part of the provenance and gives this
project no claim to ownership of that contribution.

The present prototype does not automate this track and reports no Mathlib
contribution produced by it.  It can preserve local evidence and provenance,
and it can help prepare a candidate.  Canonical status and upstream acceptance
remain external outcomes.  The
\papersectionlink{claim-faithful-publication-systems-paper.pdf}%
{systems-mathloop}{systems paper's mathematical loop} describes the authority
separation underneath this handoff.

The research graph also records objects other than theorems.  Nodes may be
endpoints, conjectures, representations, computations, proofs,
counterexamples, no-go theorems, reviews, or architecture changes.  Edges
distinguish proof dependency, implication, equivalence, refutation,
obstruction of a method, experimental support, formalisation, correction, and
publication.  The edge type matters.  A model-generated analogy can nominate
a route; it cannot become a proof dependency or a public implication without
the corresponding evidence.

As this graph grows, it may support a different kind of agent training or
retrieval: successful proofs alongside tempting arguments paired with the
exact reason they fail, and theorem statements alongside frontier
neighbourhoods showing the remaining cut.  Whether this improves
mathematical originality or merely makes models more confident within known
territory is an open empirical question.

\section{A bounded transfer hypothesis beyond mathematics}\label{sec:science}

The same abstract loop appears in experimental science: formulate a
hypothesis, design an experiment, observe an outcome, update the theory, and
retain both successful and failed interventions.  A future agent might define
a simulation or experimental script whose variables are linked to a formal
model and a semantic account of the hypothesis.  Physical observations could
then update the research graph and select the next experiment.

The mathematical implementation tests this sequence of research operations.
It supplies no evidence that the current system can conduct physics,
chemistry, or biology.  Formal verification has a special strength in mathematics because
the target objects and proof rules are digital.  In an experimental domain,
the instrument, calibration, sample, protocol, measurement uncertainty,
physical environment, and causal interpretation become independent
authorities.  A formally verified control program does not validate the
scientific hypothesis or make a laboratory safe.

Any transfer would therefore require domain scientists, controlled physical
infrastructure, ethics and safety review, material and environmental limits,
and explicit human authority over experiment selection and execution.  The
initial scope should be low-risk simulation or analysis of existing public
data.  Autonomous physical experimentation is not a present project goal.

\section{Limits and governance}\label{sec:limits}
\papersectiontarget{strategy-limits}

All eight endpoint problems remain open.  The repository contains
intermediate mathematics and deep records of failed routes, but no evidence
that a distributed run will solve any problem.  The claim that infrastructure
improves the conversion from compute to mathematics has not yet been measured
in a controlled experiment.

The prototype had no recorded external user or accepted outside contribution
by 31 August 2026. Its clone instructions, agent skills, cross-paper links,
pull-request path, and credit machinery may therefore contain friction that an
author-operated audit does not reveal. Reproducing a clean clone, reporting a
broken instruction, simplifying setup, repairing an unsafe default, or making
the contribution route easier to understand are first-class architecture
contributions even when they do not alter any mathematics.

The project remains centred on its maintainer.  The initial problem
selection, architecture, claim boundaries, review decisions, and public
presentation were all chosen through one operator-controlled process.  Open
source makes those choices inspectable and contestable; it does not make them
independent.  A serious growth phase needs external maintainers and reviewers,
published conflict and appeal rules, and a way for contributor-controlled
reports to remain visible when the project declines to adopt them.

Review capacity is limited.  More public agents may worsen the burden by
generating plausible but low-value returns.  Entry tests, evidence schemas,
automated replay, and Palomar-style selection should protect expert attention,
but automatic filters cannot award significance or community acceptance.

Contributors may dispute credit.  A role vocabulary preserves more information
than a single score, but it cannot mechanically decide authorship, priority,
or the relative weight of an idea and its proof.  The project must disclose
its policy in advance and preserve correction history.

An open corpus can be gamed.  Contributors may optimise for visible
activity, generated theorem counts, or a public model comparison.  Strong
models may have encountered repository text during training.  Security bugs
may arrive as apparently useful workflow improvements.  No automatic
leaderboard should be allowed to become the objective function.

Open access also leaves resource inequalities in place.  Compute donors,
frontier-model providers, established mathematicians, and maintainers have
different access and influence.  The strategy can lower the fixed cost of
entry and make attribution more durable.  It cannot by itself make research
participation equitable.

\section{Execution order}\label{sec:execution}

Further implementation should follow this order.

\begin{enumerate}
\item Make the README state the experiment, contributor types, exact open
      boundary, clone command, agent prompt, return paths, and credit policy.
\item Keep every problem world independently navigable from a cold clone,
      with one paper-level open section and bounded query routes.
\item Accept both mathematical and architecture returns through typed schemas,
      human review, immutable generations, and public provenance views.
\item Provide runner-neutral local instructions and safe, unprivileged checks
      before building a hosted or volunteer-compute scheduler.
\item Recruit external reviewers and maintainers before increasing agent
      throughput substantially.
\item Run controlled studies of navigation changes, model changes, and compute
      changes against immutable snapshots.
\item Only after these boundaries work should the project consider a public
      task market, donated compute service, or transfer to another scientific
      domain.
\end{enumerate}

The first three steps are partly implemented in the current public clone.
Turnkey multi-provider mining, a public volunteer-compute scheduler,
independent governance, and a cross-domain laboratory interface are not
reported capabilities.

\section{The next test}\label{sec:conclusion}

The public clone provides a route from a hard question to a bounded task,
with procedures for returning evidence and crediting accepted work.
Contributors can share the infrastructure while supplying mathematics,
compute, formalisation, software, review, or exposition.  Accepted results
extend the corpus; accepted infrastructure changes revise how later work is
produced and checked.  Both keep their provenance.

The eight Erd\H{o}s problems give this design real mathematical questions and
exact open boundaries.  Their difficulty requires the record to account for
failed approaches, while leaving room to study changes in models, compute,
infrastructure, and human direction over time.  The next test is whether
outside contributors can use the procedure, reviewers can assess their
returns, and the accumulated record improves subsequent attempts.

\appendix
\section{Public entry routes}\label{app:routes}
\papersectiontarget{strategy-entry-routes}

The public repository is \href{\repobase}{\texttt{wcook04/plectis-erdos}}.
The shortest current routes are:

\begin{center}
\small
\begin{tabular}{@{}p{0.35\linewidth}p{0.54\linewidth}@{}}
\toprule
\papertablehead
Question & Public route \\
\midrule
What is the experiment? & \repolink{README.md}{README.md} \\
Where should an agent begin? & \repolink{AGENTS.md}{AGENTS.md} \\
How can an agent explain the system? &
\repolink{skills/explain-public-system/SKILL.md}{explain-public-system skill} \\
How can an agent run the coupled research lifecycle? &
\repolink{skills/run-coupled-research-goals/SKILL.md}{coupled-goal skill} \\
How can an agent mine a frontier? &
\repolink{skills/mine-open-problem/SKILL.md}{mine-open-problem skill} \\
How is a bounded Lean change validated? &
\repolink{skills/lean-concurrent-validation/SKILL.md}{lean-concurrent-validation skill} \\
How are downstream consequences reconciled? &
\repolink{skills/propagate-research-consequences/SKILL.md}{propagate-research-consequences skill} \\
How are clone skills installed elsewhere? &
\repolink{skills/install-clone-skills/SKILL.md}{install-clone-skills skill} \\
What is proved and what remains open? & \repolink{docs/RESULTS.md}{docs/RESULTS.md} \\
Which papers and problems exist? & \repolink{docs/papers/README.md}{docs/papers/README.md} \\
How can I contribute mathematics? & \repolink{CONTRIBUTING.md}{CONTRIBUTING.md} and the
research-progress issue form \\
How can I improve the architecture? &
\repolink{docs/research-commons/ARCHITECTURE_CONTRIBUTIONS.md}{architecture contribution guide}
and the architecture-proposal issue form \\
How can an agent prepare a pull request? &
\repolink{skills/submit-pull-request/SKILL.md}{submit-pull-request skill} \\
How can I propose or add another problem? &
\repolink{skills/add-open-problem/SKILL.md}{add-open-problem skill} \\
How is a return validated? &
\repolink{skills/erdos-research-return/SKILL.md}{erdos-research-return skill} and the
\repolink{docs/research-commons/README.md}{research-commons protocol} \\
How is credit recorded? &
\repolink{docs/research-commons/CREDIT_POLICY.md}{credit policy} \\
What do Comparator and Palomar establish? &
\repolink{docs/EXTERNAL_VERIFICATION.md}{Comparator guide} and
\repolink{docs/verification/PALOMAR_QUALIFICATION.md}{Palomar qualification} \\
\bottomrule
\end{tabular}
\end{center}

\section*{Declaration of generative AI use}
Every word of this manuscript was generated by agents based on large language
models operating within Will Cook's private research system for artificial
intelligence. The formal proofs and repository software were likewise drafted
and revised by the agents through that system under Cook's direction. Cook set
the objectives and acceptance criteria, selected and reviewed the public
claims, and approved the published version. Cook assumes responsibility for
the accuracy, interpretation, and presentation of the work. Generative systems
are production tools, not authors, and supply no independent authority.

\begin{thebibliography}{99}

\bibitem{tao2026}
T. Tao, \emph{Mathematics in the age of AI}, 2026,
\href{https://arxiv.org/abs/2608.16753}{arXiv:2608.16753}.

\bibitem{mathandai2026}
Math and AI, \emph{A Severe Misalignment of AI in Mathematics},
\href{https://mathandai.org/}{declaration},
accessed 12 September 2026.

\bibitem{tao2026landscape}
T. Tao, \emph{Discussion of the difficulty landscape for mathematical
problems}, Mathstodon post, 8 September 2026,
\href{https://mathstodon.xyz/@tao/117237322160500501}{part 3 of 4},
accessed 12 September 2026.

\bibitem{boinc}
D. P. Anderson, \emph{BOINC: A Platform for Volunteer Computing}, Journal of
Grid Computing 18 (2020), 99--122,
\href{https://doi.org/10.1007/s10723-019-09497-9}{DOI}.

\bibitem{gimps}
Great Internet Mersenne Prime Search, \emph{GIMPS}, project documentation and
discovery-credit record, \href{https://www.mersenne.org/}{mersenne.org},
accessed August 2026.

\bibitem{polymath}
Polymath Project, \emph{General polymath rules},
\href{https://polymathprojects.org/general-polymath-rules/}{project rules},
accessed August 2026.

\bibitem{mathlib}
Lean community, \emph{Contributing to mathlib},
\href{https://leanprover-community.github.io/contribute/index.html}{contributor guide},
accessed August 2026.

\bibitem{carleson}
L. Becker et al., \emph{A Blueprint for the Formalization of Carleson's
Theorem on Convergence of Fourier Series}, 2025,
\href{https://arxiv.org/abs/2405.06423}{arXiv:2405.06423}.

\bibitem{erdosproblems}
T. F. Bloom, \emph{Erd\H{o}s Problems}, problem database, sources, and forum,
\href{https://www.erdosproblems.com/}{erdosproblems.com},
accessed August 2026.

\bibitem{lean4}
L. de Moura and S. Ullrich, \emph{The Lean 4 Theorem Prover and Programming
Language}, Automated Deduction---CADE 28 (2021), 625--635,
\href{https://doi.org/10.1007/978-3-030-79876-5_37}{DOI}.

\bibitem{formalconjectures}
Google DeepMind, \emph{formal-conjectures},
\href{https://github.com/google-deepmind/formal-conjectures}{software repository},
accessed August 2026.

\bibitem{palomar}
Palomar Registry, \emph{About Palomar} and \emph{Contribution policy},
\href{https://palomar-registry.org/about}{registry documentation} and
\href{https://github.com/PalomarRegistry/PalomarPolicy/blob/main/CONTRIBUTING.md}{submission standard},
accessed August 2026.

\bibitem{agenthunt}
C. E. Brown, C. Kaliszyk, and J. Urban, \emph{Agent Hunt: Bounty Based
Collaborative Autoformalization With LLM Agents}, 2026,
\href{https://arxiv.org/abs/2603.06737}{arXiv:2603.06737}.

\bibitem{leanatlas}
B. Yanahama and A. Sannai, \emph{Lean Atlas: An Integrated Proof Environment
for Scalable Human--AI Collaborative Formalization}, 2026,
\href{https://arxiv.org/abs/2604.16347}{arXiv:2604.16347}.

\bibitem{credit}
NISO, \emph{CRediT: Contributor Roles Taxonomy},
\href{https://credit.niso.org/contributor-roles-defined/}{role definitions},
accessed August 2026.

\bibitem{githubsecurity}
GitHub, \emph{Preventing pwn requests}, GitHub Actions security guidance,
\href{https://docs.github.com/en/actions/reference/security/secure-use}{documentation},
accessed August 2026.

\end{thebibliography}

\end{document}
